\chapter{Overview: Simplifying Models}
One oft necessary aspect of powder diffraction crystallography is the process of reducing the complexity of models, as a powder diffraction pattern may not have enough information to allow every parameter to be varied that could be included. 
We must introduce some reasoning to develop models that test for the things we wish to learn, but at the same time we also want to make sure that we don't use diffraction to tell us things that ain't so. 
Do not feel that powder diffraction is alone in reducing the complexity of models. 
In protein single-crystal diffraction analysis, one does not fit the position of every atom solely from the diffraction data without simplification or inclusion of some sort of \textit{ad hoc} force fields. 
We know the structures of amino acids; they reuse that knowledge. 
But when it comes to the positions of side-chains that influence how an enzyme functions, that they model most carefully. 
Even in small-molecule single-crystal crystallography, one can always expand the parameterization, using higher-order ADPs or modeling electron distributions, but there is a limit to the complexity that the data can support. 
In this set of chapters, we will look at the ways that models can be reduced in complexity using tools that GSAS-II provides, such as restraints and constraints, including rigid body constraints. 